% Source: https://artofproblemsolving.com/community/c3h2168598_bogus_equation_not_a_marathon_i_just_wanna_see_if_anyone_can_prove_it_wrong
% Stable identity: aops-topic-2168598-q01
Note that $\sqrt1=1$ throughout.

\textbf{(1) to (2) is invalid.} Line (1) says $x=2\sqrt1=2$, while line (2) says $x=2+\sqrt1=3$. Line (2) does not follow from line (1): subtracting $\sqrt1$ from both sides of (1) gives $2\sqrt1-\sqrt1=x-\sqrt1$, that is, $1=x-1$, and dividing (1) by $\sqrt1$ gives $2=\dfrac{x}{\sqrt1}$. The student has treated the product $2\sqrt1$ as if it were the sum $2+\sqrt1$.

\textbf{(2) to (3) is valid.} Squaring both sides of (2) gives $4=(x-\sqrt1)^2=x^2-2x\sqrt1+\left(\sqrt1\right)^2=x^2-2x\sqrt1+1$.

\textbf{(3) to (4) is invalid.} Adding $2x\sqrt1$ to both sides of (3) gives
\[4+2x\sqrt1=x^2+1,\]
but in (4) the term $+1$ has been lost.

\textbf{(4) to (5).} Line (2) gives $x=3$. In the faulty line (4) this gives $4+2\cdot3\cdot1=10$ on the left and $3^2=9$ on the right, which is the ``$10=9$''. In the correct line $4+2x\sqrt1=x^2+1$ it gives $10=10$, which is true.

So the two invalid steps are $\boxed{\text{from (1) to (2), and from (3) to (4)}}$. (Even with the second error repaired, the chain only shows that $x=3$ satisfies lines (2)--(4); it says nothing about line (1), whose solution is $x=2$.)
